\textbf{1. Kiến thức cần nhớ}
\begin{itemize}
\item \textbf{Bảng tần số ghép nhóm}: Dùng khi mẫu dữ liệu lớn hoặc là dữ liệu liên tục. Dữ liệu được chia thành các nửa khoảng $[a_1; a_2), [a_2; a_3), \dots, [a_k; a_{k+1})$.
\item \textbf{Giá trị đại diện} của nhóm $[a; b)$ là: $\boxed{c = \dfrac{a + b}{2}}$.
\item \textbf{Tần số tương đối ghép nhóm}: $f_i = \dfrac{m_i}{N} \times 100\%$.
\end{itemize}

\textbf{2. Phương pháp giải các dạng}

%% DANG 9D7?24-1
\textbf{Dạng 1. Lập bảng tần số ghép nhóm}
\begin{enumerate}
\item Xác định các khoảng chia $[a_i; a_{i+1})$ đều nhau theo yêu cầu đề bài.
\item Đếm số phần tử thuộc từng nửa khoảng (chú ý: lấy đầu mút trái, không lấy đầu mút phải).
\item Lập bảng gồm dòng nhóm giá trị và dòng tần số.
\end{enumerate}

%% DANG 9D7?24-2
\textbf{Dạng 2. Lập bảng tần số tương đối ghép nhóm}
\begin{enumerate}
\item Tính tỉ số phần trăm giữa tần số từng nhóm và cỡ mẫu $N$: $f_i = \dfrac{m_i}{N} \times 100\%$.
\item Lập bảng ghi nhóm và tần số tương đối.
\end{enumerate}

%% DANG 9D7?24-3
\textbf{Dạng 3. Vẽ và đọc biểu đồ tần số (tương đối) ghép nhóm}
\begin{enumerate}
\item Vẽ biểu đồ tần số ghép nhóm dạng cột liền nhau (biểu đồ Histogram): các cột liền kề nhau có đáy là các đoạn $[a_i; a_{i+1})$ trên trục hoành và chiều cao là tần số tương ứng.
\item Đọc thông tin: xác định nhóm có số lượng đông nhất, tỉ lệ học sinh đạt chuẩn\dots
\end{enumerate}
